% --- CHUNK_METADATA_START ---
% src_checksum: 8ae9592dcf848c91c59cc3e3acb4dc3792213563dee25179b017e6e0fb303aa1
% needs_review: True
% --- CHUNK_METADATA_END ---
\chapter{Independence}
% --- CHUNK_METADATA_START ---
% src_checksum: 72866aa510b3cc88f678364c0035ffe135e58e0d2ca75a0b2e5a6cc256da9ebb
% --- CHUNK_METADATA_END ---
\sloppy% --- CHUNK_METADATA_START ---
% src_checksum: e3f5b24cb90f572d00f28c5d81fd83d05f3fae30aadd742990388dbe64bb976f
% needs_review: True
% --- CHUNK_METADATA_END ---
\section{Sigma algebra generated by a random variable}% --- CHUNK_METADATA_START ---
% src_checksum: 1faa106bbadd625cd6019698b7b56a4362be86e524b9eb2e0bec9c715e8488c3
% needs_review: True
% --- CHUNK_METADATA_END ---
Let $(\Omega, \mathcal{F}, P)$ be a probability space.% --- CHUNK_METADATA_START ---
% src_checksum: 8e82de42beeff5230cfbc7b7947bdd542f3706fc3c736c6742462eff1f3f04ed
% needs_review: True
% --- CHUNK_METADATA_END ---
\begin{definition}[Sub-$\sigma$-algebra]
A subtribe $\mathcal{G}$ of $\mathcal{F}$ is a $sigma$-algebra that is included in $\mathcal{F}$, that is, $\mathcal{G} \subset \mathcal{F}$. For any $A \in \mathcal{G}$, we have $A \in \mathcal{F}$.
\end{definition}% --- CHUNK_METADATA_START ---
% src_checksum: 5a5f23c40a58d01755b12784f4ada8a93479aced98871ada8b7c3771938f2a84
% needs_review: True
% --- CHUNK_METADATA_END ---
\begin{definition}[$\sigma$-algebra of a Random Variable]
Let $X$ be a random variable (r.v.) from $(\Omega, \mathcal{F}, P)$ to $(\mathbb{R}, \mathcal{B}(\mathbb{R}))$. The sigma-algebra generated by the r.v. $X$, denoted by $\sigma(X)$, is the smallest sigma-algebra on $\Omega$ for which $X$ is measurable. It is defined as:
\[ \sigma(X) = \sigma(\{X^{-1}(B) \mid B \in \mathcal{B}(\mathbb{R})\}) \]
It should be noted that the collection of sets $\{X^{-1}(B) \mid B \in \mathcal{B}(\mathbb{R})\}$ is itself a sigma-algebra.
\end{definition}% --- CHUNK_METADATA_START ---
% src_checksum: 5f55cbd241e2d1994b7d7847f58550782734fd870a790e7ad2d2555f124a4681
% needs_review: True
% --- CHUNK_METADATA_END ---
Since $X$ is a random variable, for all $B \in \mathcal{B}(\mathbb{R})$, we have $X^{-1}(B) \in \mathcal{F}$. Therefore, $\sigma(X)$ is a sub-$sigma$-algebra of $\mathcal{F}$.% --- CHUNK_METADATA_START ---
% src_checksum: 26d37a1a0271ac64b3556f0ce7ecfc486bd56a4e54aa6b130b21b1c532f2e7cf
% needs_review: True
% --- CHUNK_METADATA_END ---
\section{Fundamental Definitions}% --- CHUNK_METADATA_START ---
% src_checksum: c503f7b81207116f9f249c533848c4033d841a3869576d48ec3f643b953852e8
% needs_review: True
% --- CHUNK_METADATA_END ---
\subsection{Independence of events}% --- CHUNK_METADATA_START ---
% src_checksum: c5580869c8b71cdfc8584d7842dde7aa9b9e4316dec1348dc0277196065d5d49
% needs_review: True
% --- CHUNK_METADATA_END ---
\begin{definition}
Let $n \ge 2$. The $n$ events $A_1, \ldots, A_n$ in $\mathcal{F}$ are called independent if, for any subset $I \subseteq \{1, \ldots, n\}$, we have:
\[ P\left(\bigcap_{i \in I} A_i\right) = \prod_{i \in I} P(A_i) \]
\end{definition}% --- CHUNK_METADATA_START ---
% src_checksum: 198f98e15af59eb08fe2a9963d4213cd208b0cc81847d9e60c12b7a8fe9514e3
% needs_review: True
% --- CHUNK_METADATA_END ---
\begin{remark}
Note that "being independent" is a much stronger condition than "being pairwise independent."
\end{remark}% --- CHUNK_METADATA_START ---
% src_checksum: 35f44f78d2e5568e527fdd3553969a7eb152ecbe1d3adb33daf214fc4c5a3811
% needs_review: True
% --- CHUNK_METADATA_END ---
\subsection{Independence of random variables}% --- CHUNK_METADATA_START ---
% src_checksum: 60b3ba1a1a15a6948f74d1fba9c3b1ad696b8dbb3932e68472ca44602a85b985
% needs_review: True
% --- CHUNK_METADATA_END ---
\begin{definition}
The $n$ random variables $X_1, \ldots, X_n$ from $\Omega$ to $\mathbb{R}$ are said to be independent if, for all Borel sets $B_1, \ldots, B_n$ of $\mathbb{R}$, we have:
\[ P(X_1 \in B_1, \ldots, X_n \in B_n) = P(X_1 \in B_1) \cdots P(X_n \in B_n) \]
where the notation $P(A, B)$ means $P(A \cap B)$.
\end{definition}% --- CHUNK_METADATA_START ---
% src_checksum: 1fcab79119a48c0baeee2dc26a76f027f14c409844d4f0116350155566e75fe5
% needs_review: True
% --- CHUNK_METADATA_END ---
\subsection{Independence of sub-$sigma$-algebras}% --- CHUNK_METADATA_START ---
% src_checksum: 3e25f722d9834e22772e5d4a7cc48af4cee78abf14e46b6308e6513e557155c0
% needs_review: True
% --- CHUNK_METADATA_END ---
\begin{definition}
The $n$ subtribes $\mathcal{F}_1, \ldots, \mathcal{F}_n$ of $\mathcal{F}$ are called independent if, for any choice of events $A_1 \in \mathcal{F}_1, \ldots, A_n \in \mathcal{F}_n$, we have:
\[ P(A_1 \cap \cdots \cap A_n) = P(A_1) \cdots P(A_n) \]
\end{definition}% --- CHUNK_METADATA_START ---
% src_checksum: 70917346416908224804ad83322da9660380a6c59aeb4fbda4318bac0e2dc727
% needs_review: True
% --- CHUNK_METADATA_END ---
\subsection{Links between the definitions}% --- CHUNK_METADATA_START ---
% src_checksum: 880d8209c64169cda0b800b501304af41069894b4497a8ac32612ec2c81e283f
% needs_review: True
% --- CHUNK_METADATA_END ---
The three previous definitions are hierarchical.% --- CHUNK_METADATA_START ---
% src_checksum: 8fecdc0ae016d746337cd485e47f36bafe5a3cfe286d0603bda6b212f7131acc
% needs_review: True
% --- CHUNK_METADATA_END ---
\begin{itemize}
    \item The events $A_1, \ldots, A_n$ are independent if and only if the indicator random variables $\mathbf{1}_{A_1}, \ldots, \mathbf{1}_{A_n}$ are independent.
    \item The random variables $X_1, \ldots, X_n$ are independent if and only if the $sigma$-algebras they generate, $\sigma(X_1), \ldots, \sigma(X_n)$, are independent.
\end{itemize}% --- CHUNK_METADATA_START ---
% src_checksum: 97bccf8e16275d46824257371aa129a191a1c88b2b2b7c5e1e433ff1f90e6d14
% needs_review: True
% --- CHUNK_METADATA_END ---
The definition based on tribes is therefore the most general. It allows handling more complex situations, such as the case of random variables taking values in arbitrary spaces.% --- CHUNK_METADATA_START ---
% src_checksum: bba51c8eacdf223fe6b7e5efb6cb61de1e98dabccc489f2761acc90b83d3839e
% needs_review: True
% --- CHUNK_METADATA_END ---
\section{General random variables}% --- CHUNK_METADATA_START ---
% src_checksum: 78e47cc510f2a412ad57ac15c0c2c1b3af0265821467b85b4dd3719ce80c3fcd
% needs_review: True
% --- CHUNK_METADATA_END ---
\begin{definition}
Let $(E, \mathcal{E})$ be a measurable space. A random variable taking values in $E$ is a function $X: \Omega \to E$ that is measurable from $(\Omega, \mathcal{F})$ to $(E, \mathcal{E})$, that is:
\[ \forall B \in \mathcal{E}, \quad X^{-1}(B) \in \mathcal{F} \]
\end{definition}% --- CHUNK_METADATA_START ---
% src_checksum: b62e18655f25cba11f24b25fd4c666d5f6eeeb7b9d8980d865ceb03253bd05f7
% needs_review: True
% --- CHUNK_METADATA_END ---
\begin{definition}[Distribution of a Random Variable]
The law of $X$ is the probability measure $P_X$ on $(E, \mathcal{E})$ defined by:
\[ \forall B \in \mathcal{E}, \quad P_X(B) = P(X \in B) = P(X^{-1}(B)) \]
\end{definition}% --- CHUNK_METADATA_START ---
% src_checksum: 2b1d31c64db00712d05830b447b0c55b700623fceafa4d9e03e78a15f8a75faa
% needs_review: True
% --- CHUNK_METADATA_END ---
This formalism makes it possible to construct random vectors and even random functions.% --- CHUNK_METADATA_START ---
% src_checksum: 24eb42ff1debc49def4101bb405f41c1bc96ea785a7abe2599064417ac953cef
% needs_review: True
% --- CHUNK_METADATA_END ---
\begin{example}
Let $X, Y$ be two random variables with values in $\mathbb{R}$. Let $Z = (X,Y)$. Then $Z$ is a random vector in $\mathbb{R}^2$. We say that the law of $Z$ admits a density if there exists a measurable function $f: \mathbb{R}^2 \to \mathbb{R}_+$ such that for all $B \in \mathcal{B}(\mathbb{R}^2)$:
\[ P_Z(B) = P(Z \in B) = \iint_B f(x,y) \,dx\,dy \]
\end{example}% --- CHUNK_METADATA_START ---
% src_checksum: 9271cb7837f59f5e1c25b7c02e7e66ee941547f5f2eb92fb88bac9454eb5044b
% needs_review: True
% --- CHUNK_METADATA_END ---
\section{Product Laws}
% --- CHUNK_METADATA_START ---
% src_checksum: 2560fcc5d875042016e2326dd753542eb7e9ac096d0737d49ff0ad949552d8e4
% needs_review: True
% --- CHUNK_METADATA_END ---
Let $(E_1, \mathcal{E}_1, \mu_1), \ldots, (E_n, \mathcal{E}_n, \mu_n)$ $n$ be probability spaces.
A natural question arises: does there exist a probability space $(\Omega, \mathcal{F}, P)$ on which are defined $n$ random variables $X_1, \ldots, X_n$ such that:
% --- CHUNK_METADATA_START ---
% src_checksum: bb23914095a9bac07893cbf743ef6b598e4157e7d540fcb6d2b0e14dae2858d8
% needs_review: True
% --- CHUNK_METADATA_END ---
\begin{itemize}
    \item $X_1, \ldots, X_n$ are independent.
    \item For all $i \in \{1, \ldots, n\}$, $X_i$ takes values in $E_i$ and has the distribution $\mu_i$.
\end{itemize}
% --- CHUNK_METADATA_START ---
% src_checksum: 48775d168978c992c2137be4fa1a4cdb1cbb0093c88d7e8a501afa42280cddee
% needs_review: True
% --- CHUNK_METADATA_END ---
The answer is affirmative and is provided by the construction of the product measure.
% --- CHUNK_METADATA_START ---
% src_checksum: e88c9c2c14995582ca1203774c7dca5efca0874d10ed19673fca7a904332a13b
% needs_review: True
% --- CHUNK_METADATA_END ---
\subsection{Construction of independent variables}
% --- CHUNK_METADATA_START ---
% src_checksum: 1ec7793e32187fc216e08a09dd9ede462c39b30c15c92e93dd2c05030fdcb58c
% needs_review: True
% --- CHUNK_METADATA_END ---
Consider the Cartesian product space $E = E_1 \times \cdots \times E_n$. We equip it with the product sigma-algebra $\mathcal{E} = \mathcal{E}_1 \otimes \cdots \otimes \mathcal{E}_n$, which is the sigma-algebra generated by the "measurable rectangles" of the form $B_1 \times \cdots \times B_n$, where $B_i \in \mathcal{E}_i$ for each $i$.% --- CHUNK_METADATA_START ---
% src_checksum: a250c4575ecae3f9b7563896c6362bed10f7f169cf4ce14abe680f84b11892cf
% needs_review: True
% --- CHUNK_METADATA_END ---
\begin{theorem}[Product Law]
There exists a unique probability measure $\mu$ on the measurable space $(E, \mathcal{E})$ such that for all $B_1 \in \mathcal{E}_1, \ldots, B_n \in \mathcal{E}_n$:
\[ \mu(B_1 \times \cdots \times B_n) = \mu_1(B_1) \cdots \mu_n(B_n) \]
This measure is called the product measure of $\mu_1, \ldots, \mu_n$ and is denoted by $\mu = \mu_1 \otimes \cdots \otimes \mu_n$.
\end{theorem}% --- CHUNK_METADATA_START ---
% src_checksum: 231f4ad8ecb7d9da3a4498db4e8187299e7a06e4be4b177a7fb0ef5e91fcc9e1
% needs_review: True
% --- CHUNK_METADATA_END ---
To answer our initial question, we simply need to ask:% --- CHUNK_METADATA_START ---
% src_checksum: 29ef2c80a0980a4c789c91c84f30059b199ac2a128123b7f8f9947874c8edb17
% --- CHUNK_METADATA_END ---
\begin{itemize}
    \item $\Omega = E = E_1 \times \cdots \times E_n$
    \item $\mathcal{F} = \mathcal{E} = \mathcal{E}_1 \otimes \cdots \otimes \mathcal{E}_n$
    \item $P = \mu = \mu_1 \otimes \cdots \otimes \mu_n$
\end{itemize}% --- CHUNK_METADATA_START ---
% src_checksum: 9285da32a1181c28a497b210c1daf9889111effde545639a3bfc88f134e5f6fb
% needs_review: True
% --- CHUNK_METADATA_END ---
And to define the random variables $X_i$ as the projections onto the $i$-th component, for $i \in \{1, \ldots, n\}$:
\[ X_i: \Omega \to E_i, \quad \omega = (\omega_1, \ldots, \omega_n) \mapsto \omega_i \]% --- CHUNK_METADATA_START ---
% src_checksum: f6421c426b0418f007ba95d78c50dac318fb0f593d72802c64bc30fd5e61900f
% needs_review: True
% --- CHUNK_METADATA_END ---
\begin{proof}
Let us verify that these $X_i$ have the desired properties.
\begin{enumerate}
    \item \textbf{Law of $X_i$: } Let $i \in \{1, \ldots, n\}$ and $B_i \in \mathcal{E}_i$. The law of $X_i$ is given by:
    \begin{align*}
    P_{X_i}(B_i) &= P(X_i \in B_i) = \mu(\{\omega = (\omega_1, \ldots, \omega_n) \in \Omega \mid \omega_i \in B_i\}) \\
    &= \mu(E_1 \times \cdots \times E_{i-1} \times B_i \times E_{i+1} \times \cdots \times E_n) \\
    &= \mu_1(E_1) \cdots \mu_{i-1}(E_{i-1}) \mu_i(B_i) \mu_{i+1}(E_{i+1}) \cdots \mu_n(E_n)
    \end{align*}
    Since $\mu_j(E_j)=1$ for all $j$, we obtain $P_{X_i}(B_i) = \mu_i(B_i)$. Thus, $X_i$ indeed has the law $\mu_i$.
    \item \textbf{Independence of the $X_i$: } Let $B_1 \in \mathcal{E}_1, \ldots, B_n \in \mathcal{E}_n$. By definition of independence:
    \begin{align*}
    P(X_1 \in B_1, \ldots, X_n \in B_n) &= \mu(\{\omega = (\omega_1, \ldots, \omega_n) \mid \forall i, \omega_i \in B_i\}) \\
    &= \mu(B_1 \times \cdots \times B_n) \\
    &= \mu_1(B_1) \cdots \mu_n(B_n) \\
    &= P(X_1 \in B_1) \cdots P(X_n \in B_n)
    \end{align*}
    The variables $X_1, \ldots, X_n$ are therefore independent.
\end{enumerate}
\end{proof}% --- CHUNK_METADATA_START ---
% src_checksum: b23599f69402742f0ccccf504b71286fe3727b971d726ff6eafd6eeb9dedd0c1
% needs_review: True
% --- CHUNK_METADATA_END ---
\subsection{Application: i.i.d. variables}% --- CHUNK_METADATA_START ---
% src_checksum: df846b5b43d63cf800537aa41ea5f60425543d916fc0d25dc33877e506a57ec1
% needs_review: True
% --- CHUNK_METADATA_END ---
A particular important case is where $(E_i, \mathcal{E}_i) = (\mathbb{R}, \mathcal{B}(\mathbb{R}))$ for every $i$, and where all laws are identical: $\mu_1 = \cdots = \mu_n = \nu$. The theorem then guarantees the existence of a probability space in which are defined $n$ random variables $X_1, \ldots, X_n$ \textbf{independent and identically distributed} (i.i.d.), with common law $\nu$.% --- CHUNK_METADATA_START ---
% src_checksum: 5a5c06a158e702c0b311ead5b11067baf8cef691b077b22c5fe5960341ecde91
% needs_review: True
% --- CHUNK_METADATA_END ---
\section{Grouping by blocks}% --- CHUNK_METADATA_START ---
% src_checksum: 4237dbb9be4b2be003e7922dfc0f94620d196eedbba9c628ee35cd619bb50465
% needs_review: True
% --- CHUNK_METADATA_END ---
Independence is maintained through grouping.% --- CHUNK_METADATA_START ---
% src_checksum: 61d1a35167ede03e9dd92741db894a57dac2587972af88b1c98af2de8779b747
% needs_review: True
% --- CHUNK_METADATA_END ---
\begin{theorem}
Let $\mathcal{F}_1, \ldots, \mathcal{F}_n$ be independent sub-$sigma$-algebras. Let $I$ and $J$ be two disjoint subsets of $\{1, \ldots, n\}$. Then the $sigma$-algebras $\mathcal{G} = \sigma(\bigcup_{i \in I} \mathcal{F}_i)$ and $\mathcal{H} = \sigma(\bigcup_{j \in J} \mathcal{F}_j)$ are independent.
\end{theorem}% --- CHUNK_METADATA_START ---
% src_checksum: 0527fed64ca0ca797de1e7d92439d53ebea1aa163aebcbf61e71b27e3ac1e6f7
% needs_review: True
% --- CHUNK_METADATA_END ---
The proof of this theorem relies on the uniqueness extension lemma, also known as the monotone class lemma or Dynkin's $\pi-\lambda$ theorem.% --- CHUNK_METADATA_START ---
% src_checksum: 018bf99f108922701b94c12523d10ce51122c1c3bc3224e8b81aee7cf8a6136e
% needs_review: True
% --- CHUNK_METADATA_END ---
\begin{lemma}[Uniqueness of the Extension]
Let $P_1$ and $P_2$ be two probability measures on a measurable space $(\Omega, \mathcal{F})$. Suppose there exists a collection $\mathcal{C}$ of events such that:
\begin{enumerate}
    \item[i)] $\mathcal{C}$ generates $\mathcal{F}$, i.e., $\sigma(\mathcal{C}) = \mathcal{F}$.
    \item[ii)] $\mathcal{C}$ is stable under finite intersection (it is a $\pi$-system), i.e., $\forall A, B \in \mathcal{C}, A \cap B \in \mathcal{C}$.
    \item[iii)] $P_1$ and $P_2$ coincide on $\mathcal{C}$, i.e., $\forall C \in \mathcal{C}, P_1(C) = P_2(C)$.
\end{enumerate}
Then $P_1 = P_2$ (they coincide on the entire $\sigma$-algebra $\mathcal{F}$).
\end{lemma}% --- CHUNK_METADATA_START ---
% src_checksum: ab31d60b5d1fbf7027982e1c042518950a62a6e6b83c030bd3a457d4bf541d0b
% not-translated-due-to-exception: True
% --- CHUNK_METADATA_END ---
\begin{proof}[Proof Idea of the Theorem]
On veut montrer que $\forall D \in \mathcal{G}, \forall E \in \mathcal{H}, P(D \cap E) = P(D)P(E)$.
\begin{enumerate}
    \item On fixe d'abord un événement $D$ dans une classe simple génératrice de $\mathcal{G}$, par exemple $D = \bigcap_{i \in I} D_i$ avec $D_i \in \mathcal{F}_i$. On suppose $P(D)>0$.
    \item On définit deux mesures de probabilité sur $(\Omega, \mathcal{H})$ : $P_1(E) = P(E|D) = \frac{P(E \cap D)}{P(D)}$ et $P_2(E) = P(E)$.
    \item On montre que $P_1$ et $P_2$ coïncident sur une collection $\mathcal{C}$ qui est un $\pi$-système engendrant $\mathcal{H}$. La collection des "rectangles" $E = \bigcap_{j \in J} E_j$ avec $E_j \in \mathcal{F}_j$ convient. En effet, pour un tel $E$ :
    \begin{align*}
    P(E \cap D) &= P\left( \left(\bigcap_{j \in J} E_j\right) \cap \left(\bigcap_{i \in I} D_i\right) \right) \\
    &= \left( \prod_{j \in J} P(E_j) \right) \left( \prod_{i \in I} P(D_i) \right) \quad \text{(par indépendance des } \mathcal{F}_k) \\
    &= P(E) P(D)
    \end{align*}
    Donc $P_1(E) = P(E|D) = P(E) = P_2(E)$.
    \item Par le lemme d'unicité, $P_1(E) = P_2(E)$ pour tout $E \in \mathcal{H}$, ce qui signifie $P(E \cap D) = P(E)P(D)$ pour tout $E \in \mathcal{H}$ et pour $D$ de la forme "rectangle".
    \item On refait un raisonnement similaire en fixant $E \in \mathcal{H}$ et en faisant varier $D$ dans $\mathcal{G}$.
\end{enumerate}
\end{proof}% --- CHUNK_METADATA_START ---
% src_checksum: 5eb1cfd6dad53a3c4f9388d6aa441017d6b821287407e41a17e04c84edcf5ab7
% needs_review: True
% --- CHUNK_METADATA_END ---
This theorem is generalized by recursion.% --- CHUNK_METADATA_START ---
% src_checksum: 08b2a2e1ee0421128749690036717a687ca86797a1f365ae6bb54518ab6a4cc7
% needs_review: True
% --- CHUNK_METADATA_END ---

\begin{proposition}
Let $\mathcal{F}_1, \ldots, \mathcal{F}_n$ be independent $sigma$-algebras. Let $I_1, \ldots, I_k$ be subsets of $\{1, \ldots, n\}$ that are pairwise disjoint. Then the $sigma$-algebras $\mathcal{G}_j = \sigma(\bigcup_{i \in I_j} \mathcal{F}_i)$ for $j=1, \ldots, k$ are independent.
\end{proposition}
% --- CHUNK_METADATA_START ---
% src_checksum: d474c2633764831730d8205444fe3e34c473067423d1f96da6eea30fdce6edb9
% needs_review: True
% --- CHUNK_METADATA_END ---
A direct corollary for random variables is as follows:% --- CHUNK_METADATA_START ---
% src_checksum: f16316b7951c1e6ce0ed525af9e33abf8b3c8c37593add92b03b8ee7dacde29f
% needs_review: True
% --- CHUNK_METADATA_END ---
\begin{theorem}[Corollary]
Let $X_1, \ldots, X_n$ be independent random variables. Let $I_1, \ldots, I_k$ be disjoint subsets of $\{1, \ldots, n\}$. The random vectors $V_j = (X_i)_{i \in I_j}$ for $j=1, \ldots, k$ are independent.
\end{theorem}% --- CHUNK_METADATA_START ---
% src_checksum: 49d27f6e1c55c6beccffc943895c53715b2ec6523772a488de0e6520614a791f
% needs_review: True
% --- CHUNK_METADATA_END ---
\section{Expectation and independence}% --- CHUNK_METADATA_START ---
% src_checksum: dadeccece1d41a89eb4b71712f5bf022ac6c76d3bf6b1a508b4de0d0898e4f87
% needs_review: True
% --- CHUNK_METADATA_END ---
Independence has a very important consequence for the calculation of expectations.% --- CHUNK_METADATA_START ---
% src_checksum: 46d2bfe8b8ba1d0abfa954a029f30fd74a290387761d804cfa83c49481b80574
% needs_review: True
% --- CHUNK_METADATA_END ---
\begin{proposition}
Let $X$ and $Y$ be two real random variables. The following two conditions are equivalent:
\begin{enumerate}
    \item[i)] $X$ and $Y$ are independent.
    \item[ii)] For all Borel functions $f, g: \mathbb{R} \to \mathbb{R}$ such that $f(X)$, $g(Y)$, and $f(X)g(Y)$ are integrable, we have:
    \[ E[f(X)g(Y)] = E[f(X)]E[g(Y)] \]
\end{enumerate}
\end{proposition}% --- CHUNK_METADATA_START ---
% src_checksum: eeb51de311e9d66924621bbc307a28dfc2cc8768ce17c8e70dafd71d493a7e12
% needs_review: True
% --- CHUNK_METADATA_END ---
\begin{proof}
\textbf{(ii) $\implies$ (i) :} We just need to choose indicator functions. Let $A, B \in \mathcal{B}(\mathbb{R})$. We set $f = \mathbf{1}_A$ and $g = \mathbf{1}_B$. These functions are Borel measurable and bounded, so the random variables $\mathbf{1}_A(X)$, $\mathbf{1}_B(Y)$ and their product are integrable.
The equality becomes:
\[ E[\mathbf{1}_A(X) \mathbf{1}_B(Y)] = E[\mathbf{1}_A(X)] E[\mathbf{1}_B(Y)] \]
Which translates to:
\[ P(X \in A, Y \in B) = P(X \in A) P(Y \in B) \]
Since this is true for all Borel sets $A$ and $B$, $X$ and $Y$ are independent.
\textbf{(i) $\implies$ (ii) :} The proof proceeds by successive steps (monotone class theorem).
\begin{enumerate}
    \item \textbf{Indicator functions :} If $f = \mathbf{1}_A$ and $g = \mathbf{1}_B$, the equality is the definition of independence, as seen above. It is therefore true.
    \item \textbf{Positive simple functions :} By linearity of expectation, the equality extends to positive simple functions (finite linear combinations with positive coefficients of indicator functions).
    \item \textbf{Positive measurable functions :} Any positive measurable function can be approximated by a non-decreasing sequence of positive simple functions. By the monotone convergence theorem, the equality holds in the limit.
    \item \textbf{General integrable functions :} We decompose $f = f^+ - f^-$ and $g = g^+ - g^-$, where $f^+, f^-, g^+, g^-$ are positive functions. We apply the previous result to each pair and use the linearity of expectation.
    \begin{align*}
    E[f(X)g(Y)] &= E[(f^+(X) - f^-(X))(g^+(Y) - g^-(Y))] \\
    &= E[f^+(X)g^+(Y) - f^+(X)g^-(Y) - f^-(X)g^+(Y) + f^-(X)g^-(Y)] \\
    &= E[f^+(X)g^+(Y)] - E[f^+(X)g^-(Y)] - E[f^-(X)g^+(Y)] + E[f^-(X)g^-(Y)] \\
    &= E[f^+(X)]E[g^+(Y)] - E[f^+(X)]E[g^-(Y)] - E[f^-(X)]E[g^+(Y)] + E[f^-(X)]E[g^-(Y)] \\
    &= (E[f^+(X)] - E[f^-(X)])(E[g^+(Y)] - E[g^-(Y)]) \\
    &= E[f(X)]E[g(Y)]
    \end{align*}
\end{enumerate}
\end{proof}% --- CHUNK_METADATA_START ---
% src_checksum: d2c8c5f6c70b3848700702c73e0bcb1cee094da6ce4dc5f9fbf5193899333faf
% needs_review: True
% --- CHUNK_METADATA_END ---
An immediate and fundamental corollary:% --- CHUNK_METADATA_START ---
% src_checksum: cd70f1f94461c4e7f8332a8f774047fb5f1c26f1938d048cfb514578c36ec61e
% needs_review: True
% --- CHUNK_METADATA_END ---
\begin{theorem}[Corollary]
Let $X$ and $Y$ be two independent random variables such that $E[|X|] < \infty$ and $E[|Y|] < \infty$. Then $XY$ is integrable ($E[|XY|] < \infty$) and we have:
\[ E[XY] = E[X]E[Y] \]
\end{theorem}% --- CHUNK_METADATA_START ---
% src_checksum: b7e7a217c89a9594c47ab59aa7bd4c8fb65adae8ab7d3eb2fba4a8bdf8d0af20
% needs_review: True
% --- CHUNK_METADATA_END ---
\section{Covariance}% --- CHUNK_METADATA_START ---
% src_checksum: f1504fc6ea92c4428e828f974127466429ddc3cf759dfaf5737b331d345565a0
% needs_review: True
% --- CHUNK_METADATA_END ---

\begin{definition}
Let $X$ and $Y$ be two random variables on $(\Omega, \mathcal{F}, P)$ that are integrable, and whose product $XY$ is also integrable.
The covariance of $X$ and $Y$ is defined by:
\[ \text{cov}(X,Y) = E[(X-E[X])(Y-E[Y])] \]
\end{definition}
% --- CHUNK_METADATA_START ---
% src_checksum: 153a54d565b154d183f76b6e496c99c203d70233e60f90727867dc170f26a6de
% needs_review: True
% --- CHUNK_METADATA_END ---
By developing the expression, we obtain a more practical formula:% --- CHUNK_METADATA_START ---
% src_checksum: 392e0b526e518d1ce9fed2ca30d711c56193ade4d000e9621b4b52945c20c881
% needs_review: True
% --- CHUNK_METADATA_END ---
\begin{align*}
\text{cov}(X,Y) &= E[XY - X E[Y] - Y E[X] + E[X]E[Y]] \\
&= E[XY] - E[X]E[Y] - E[Y]E[X] + E[X]E[Y] \\
&= E[XY] - E[X]E[Y]
\end{align*}% --- CHUNK_METADATA_START ---
% src_checksum: 96dd4d37c38c44e18a618bd4abe29a3594a5d6b90d2cc22f0a38da0dd9d242ed
% needs_review: True
% --- CHUNK_METADATA_END ---
Covariance is a measure of the linear dependence between two variables.% --- CHUNK_METADATA_START ---
% src_checksum: 62dfb52974f8f4b69b64b92c6124ae7bc4b0f65f96e5c2ce869356cb49223edb
% needs_review: True
% --- CHUNK_METADATA_END ---
\begin{proposition}
If $X$ and $Y$ are independent and integrable, then their covariance is zero:
\[ \text{cov}(X,Y) = 0 \]
\end{proposition}% --- CHUNK_METADATA_START ---
% src_checksum: ece2b213176544ef72457013d69de400844a481fb3ff7b73a37ab13c4acb5e9f
% needs_review: True
% --- CHUNK_METADATA_END ---
\begin{proof}
If $X$ and $Y$ are independent and integrable, by the previous corollary, $E[XY] = E[X]E[Y]$. Consequently, $\text{cov}(X,Y) = E[XY] - E[X]E[Y] = 0$.
\end{proof}% --- CHUNK_METADATA_START ---
% src_checksum: 0b946742763cb59d64e3ece643a0b4bb71b797ced1aa8dfb0e84163e4639650e
% needs_review: True
% --- CHUNK_METADATA_END ---
\begin{remark}
The converse is false! Two variables may have zero covariance (they are said to be uncorrelated) without being independent. Covariance only captures linear dependencies.
\end{remark}
